\title{The Shift to Agentic AI: Evidence from Codex}
\begin{document}
\maketitle

\begin{abstract}
THE SHIFT TO AGENTIC AI: EVIDENCE FROM CODEX Drew Johnston1,* David Holtz2,1 Alex Martin Richmond1 Christopher Ong1 Prasanna Tambe3,1 Aaron Chatterji1,4 1OpenAI 2Columbia Business School 3University of Pennsylvania, Wharton School 4Duke University Fuqua School of Business We analyze usage data from OpenAI’s Codex tool to present large-scale evidence of how agentic AI technology, which can take actions on a user’s behalf, changes how people work. We use an automated, privacy-protecting pipeline to contrast usage across three populations: external personal-account users, external organizational-account users, and workers within OpenAI. We find that agentic AI usage is growing rapidly: the number of active users has grown more than fivefold in the first half of 2026, with the most rapid increase occurring outside the initial audience of software developers. Uptake is uneven across contexts: within OpenAI, Codex usage is nearly universal and has largely replaced business usage of ChatGPT. We document a similar shift to agentic tooling outside Ope- nAI, particularly within organizations, although external adoption remains lower and more uneven. In addition to headline usage figures, we observe measures of sophistica- tion, and find that a growing number of users have used Codex to change their workflows substantially. We find that more than 10\% of users manage three or more concurrent Codex agents at some point each week and that 26.6\% use skills, which allow users to share in- st
\end{abstract}

\section{Recovered Text}
THE SHIFT TO AGENTIC AI: EVIDENCE
FROM CODEX
Drew Johnston1,*
David Holtz2,1
Alex Martin Richmond1
Christopher Ong1
Prasanna Tambe3,1
Aaron Chatterji1,4
1OpenAI
2Columbia Business School
3University of Pennsylvania, Wharton School
4Duke University Fuqua School of Business
We analyze usage data from OpenAI’s Codex tool to present large-scale evidence of how
agentic AI technology, which can take actions on a user’s behalf, changes how people
work. We use an automated, privacy-protecting pipeline to contrast usage across three
populations: external personal-account users, external organizational-account users, and
workers within OpenAI. We find that agentic AI usage is growing rapidly: the number
of active users has grown more than fivefold in the first half of 2026, with the most rapid
increase occurring outside the initial audience of software developers. Uptake is uneven
across contexts: within OpenAI, Codex usage is nearly universal and has largely replaced
business usage of ChatGPT. We document a similar shift to agentic tooling outside Ope-
nAI, particularly within organizations, although external adoption remains lower and
more uneven. In addition to headline usage figures, we observe measures of sophistica-
tion, and find that a growing number of users have used Codex to change their workflows
substantially. We find that more than 10\% of users manage three or more concurrent Codex
agents at some point each week and that 26.6\% use skills, which allow users to share in-
structions for complex workflows. Alongside these changes in usage practices, request
complexity has increased: since the start of the year, the share of individual Codex users
who submit at least one request for a task estimated to require more than eight hours for
an experienced human to complete has increased nearly tenfold. Concurrently, output has
grown rapidly—in June 2026, the median OpenAI employee in a legal role generated 13
times more monthly output tokens across Codex and ChatGPT than they did in Novem-
ber 2025, while the median researcher generated more than 50 times as many. We conclude
by discussing the implications of these patterns for productivity, job reorganization, and
workforce restructuring.
The authors are grateful to Cassandra Duchan Solis for her substantial contributions to this paper, and to Tejal Patwardhan,
Suproteem Sarkar, and Katy Shi for helpful early conversations and feedback. Andrew Hillis and Kevin Liu provided help-
ful insights about the data infrastructure underlying many of the analyses in this paper. Gawesha Weeratunga arranged the
data access supporting our analyses of usage within organizations. Daniella Eguiguren and Kirsty Le produced the web vi-
sualization, which was developed by Matt Nichols, Shawn O’Mara, and Dion Peters. Charlene Tsang designed portions of
the visualization. Rachel Brown, Laurance Fauconnet, Rob Friedlander, and Laurie Jones provided editorial feedback. Laura
Bisesto, Sarah Friar, Zak Weiler, and Meng Jia Yang reviewed this paper. David Holtz and Prasanna Tambe contributed to this
work in their capacity as paid contractors for OpenAI.
*Corresponding author: drewjohnston@openai.com
1
1
Introduction
Generative AI systems increasingly differ in the extent to which they can act on a user’s behalf, or in
other words, the extent to which they are agentic. Traditional chatbot interfaces are primarily conver-
sational: users ask questions or provide instructions, and the model generates responses. Agentic AI
tools extend this paradigm by enabling users to delegate multi-step tasks to AI systems that can au-
tonomously use external tools, inspect files, execute commands, and create or modify artifacts.
This paper studies this shift to agentic AI by analyzing data on the adoption and use of Codex, Ope-
nAI’s agentic1 coding and work platform, which was initially released in April 2025.2 Codex was origi-
nally designed for software development, a domain where AI outputs are useful, economically impor-
tant, and comparatively easy to verify. However, the usage we study extends beyond software: Codex is
used for drafting documents, analyzing data, coordinating communication, and other knowledge-work
tasks outside coding.
To study how agentic AI use diffuses across settings and expands across task domains, we analyze
Codex usage data. These data allow us to document adoption, task allocation, and depth of use, using
automated classifiers to extract aggregated and anonymized insights about AI usage without researchers
reading the underlying messages themselves.
We conduct this analysis across three distinct populations: individual users, organizational users,
and OpenAI workers.3 Individual users provide a broad-market view of early adoption, where use is
heterogeneous and workflow integration is often limited. Organizational users show how Codex is used
in business settings outside OpenAI, where adoption may depend on local workflows, permissions, se-
curity requirements, and user familiarity. OpenAI workers provide a unique view into what usage is at
the frontier. OpenAI is an unusually favorable environment for agentic AI: workers are highly familiar
with frontier models, usage is cheap at the margin, organizational buy-in is high, training and informal
knowledge sharing are common, and many workflows are close to the systems being developed.4 Ope-
nAI usage is therefore not representative of the typical organization today. However, it provides a view
of what agentic AI use may look like in the future, when adoption frictions are minimal.
We organize the evidence around four stylized facts.
First, the shift to agentic AI is rapid but uneven. Over the first six months of 2026, weekly active
Codex usage increased sharply, yet agentic tooling remains much less broadly used than ChatGPT. The
shift is smallest among individual users, larger among organizational users, and largest among OpenAI
1In this paper, for brevity, we refer to ChatGPT as a conversational AI tool and Codex as an agentic AI tool. In practice, the dis-
tinction is not absolute. ChatGPT includes agentic capabilities, such as web browsing and code execution, while some Codex
interactions are purely conversational. In the week preceding June 11, 2026, 60.3\% of Codex turns and 21.9\% of ChatGPT
turns invoked at least one external tool. Tool use is an imperfect proxy for agency: some tool invocations are part of simple
conversational interactions, while some agentic workflows involve limited tool use.
2This initial release was limited in scope and available only via a command line interface. Later releases broadened the scope
of users eligible to use Codex and the variety of surfaces through which it could be accessed.
3We refer to users on personal plans, including Free, Go, Plus, and Pro, as Individual users. We refer to users on Business and
Enterprise plans as Organizational users.
4Internally at OpenAI, there are no quantity restrictions on AI usage, and there is significant internal knowledge sharing about
AI capabilities, skills, and usage patterns.
2
workers. The contrast is sharpest when measured through the use of output tokens rather than numbers
of active users.5 Among OpenAI workers, Codex largely replaced ChatGPT as the main interface for
work-related AI use: as of June 11, 2026, Codex accounts for 99.8\% of output tokens these workers
generate across Codex and ChatGPT. Among organizational users, the corresponding share is 63.3\%,
while among individual users it is 16.5\%. Adoption also varies by job role. Technical roles adopt earlier,
but non-developer use has grown quickly, especially within OpenAI.
Second, Codex use is strongly oriented toward delegated production. Users ask the model to carry
out concrete work tasks like debugging, refactoring, validating changes, configuring applications, draft-
ing documents, and analyzing data. These activities are better understood as production than as consul-
tation: users are asking Codex to do work, not only to provide advice or information.6 This distinction
is central to the shift we study. In an agentic interface, usage can correspond to a delegated workflow
rather than a single conversational exchange. We find that the complexity of tasks users request has
increased over time, with an increasing share of users delegating tasks that would take an experienced
human more than a full day of work to complete.
Third, Codex use is anchored in software production, but is broader in contexts in which adoption
is deepest. Across user populations, the largest share of tasks is closely tied to software work, including
code implementation, code understanding, code validation, engineering operations, and application
management. But users delegate more than code generation, using Codex across the broader software-
production life cycle: they also use it to understand existing systems, configure environments, validate
changes, manage repositories, and produce documentation. Among OpenAI workers, usage extends
further into research, planning, communication, data analysis, product work, recruiting, sales, and other
non-engineering activities.
Fourth, intensive users organize Codex use around large, repeatable, and parallel workflows. We
measure this in several ways. Output volume measures the amount of AI-mediated work users gener-
ate. Task-complexity measures the estimated time it would take an experienced human to complete the
tasks that users delegate. Skill use measures whether users invoke reusable instructions, capabilities,
or tool integrations. Runtime measures how much active agent work occurs on a user’s behalf. Con-
currency measures whether users run multiple Codex turns at the same time. Across these measures,
heavy users look qualitatively different from occasional users. They are more likely to use skills, run
longer and more complex tasks, and operate multiple agents concurrently. OpenAI workers provide the
clearest view of this pattern: among the most intensive users within OpenAI, Codex is less an assistant
answering requests and more like a workflow system in which the user delegates, monitors, reviews,
and coordinates multiple streams of work.
The evidence illustrates why agentic AI is more than a more capable version of conversational AI:
users can hand off larger pieces of work to systems that inspect context, use tools, and modify arti-
5Tokens are the basic unit of information read and generated by most AI models.
6This contrasts with conversational AI usage patterns documented in Chatterji et al. (2025). In that context, researchers found
that a larger share of work focused on “asking” (which represented nearly half of all prompts), than on “doing”.
3
facts. This distinction changes what adoption means. For agentic AI, the important margins are not only
whether a person uses the tool, but what work they delegate, how much execution the system performs,
and whether users begin to organize workflows around repeated or parallel delegation.
The rest of the paper proceeds as follows. Section 2 describes the related literature and connects it
to the work of this paper. Section 3 studies who is using agentic tooling, documenting growth in Codex
usage across users in individual, organizational, and internal OpenAI accounts. Within these categories,
we study differences in usage across personas, job functions, and seniority levels. Section 4 examines
what people use Codex for, comparing task composition across user populations and worker groups.
Section 5 studies how people use Codex once they adopt it, focusing on output volume, repeatable
workflows, task complexity, runtime, and concurrency. The conclusion summarizes what these usage
patterns show about the frontier of agentic AI use.
2
Related literature
In studying the diffusion and use of Codex, we connect the emerging literature on agentic AI to three
related literatures that study generative AI adoption, delegation and verification in human–AI systems,
and the organizational complements that shape technology diffusion.
First, we extend work on the adoption and use of generative AI. Prior work measures occupational
exposure to large language models, generative AI adoption by workers, and usage of conversational AI
systems (Eloundou et al. 2024; Felten et al. 2023; Bick et al. 2026; Chatterji et al. 2025; Handa et al. 2025).
This literature documents the rapid diffusion of generative AI across occupations and industries while
highlighting substantial heterogeneity in adoption and use. Recent studies also describe how individu-
als employ conversational AI systems in practice, finding that information retrieval, learning, writing,
content creation, and practical guidance account for a large share of observed usage (Chatterji et al.
2025).
We extend this work by studying agentic AI tools, where the relevant unit of analysis is a delegated
workflow rather than a conversation. Like other papers in this literature, our goal is to understand the
underlying tasks people perform with new technologies by analyzing interactions between users and AI
systems. However, the agentic nature of Codex allows us to move beyond understanding conversations
and to directly observe actions performed on a user’s behalf.
This distinction connects our work to an emerging literature on agentic AI tools. The closest pa-
pers to our work study usage of Perplexity’s agentic products. Yang et al. (2025) study the adoption
of a general-purpose browser agent and show that early agentic use is broad, spanning productivity,
learning, shopping, media, career, and travel tasks, with usage varying across personal, professional,
and educational contexts. Yang et al. (2026) compare Perplexity’s Search and Computer products us-
ing matched tasks, showing that the more autonomous Computer product performs substantially more
work on users’ behalf, reduces completion time, and expands the scope of tasks users attempt. A related
software-focused literature studies coding agents in practice. Sarkar (2026) uses data from an AI-assisted
programming platform to show that agents shift software work away from implementation activities
4
such as typing code and toward delegation. Baumann et al. (2026) introduce SWE-chat, a dataset of
real coding-agent sessions from open-source developers, and show that coding-agent use in the wild
is heterogeneous, involves substantial user correction, and often differs from benchmark-style evalua-
tions. More broadly, a growing literature examines how agentic systems reshape the allocation of work
between humans and machines, creating new opportunities for delegation while introducing new re-
quirements for supervision, verification, and coordination.
These questions connect the emerging literature on agentic AI to a long tradition of research on
technology diffusion, organizational complements, and workplace change. A central insight from this
literature is that productivity gains from new technologies depend on complementary investments in
business processes, worker skills, organizational design, and intangible capital. For example, Brynjolfs-
son et al. (2019) argue that the gap between rapid AI progress and slow measured productivity growth
reflects the time required for organizations to develop complementary innovations and organizational
changes.
A shift toward agentic AI changes not only what tasks can be automated but also how work is
organized. Hitzig et al. (2026) argue that agentic systems move human–AI interaction from assistance
toward delegation, making supervision, verification, and coordination central determinants of value cre-
ation while increasing returns to domain expertise. More broadly, realizing value from agentic AI may
depend on the diffusion of AI capabilities throughout the workforce and the growing importance of
domain expertise in supervising and coordinating delegated work (Tambe 2026; Hitzig et al. 2026). Con-
sistent with this view, Demirer et al. (2026b) show that large task-level productivity gains may translate
only imperfectly into final output because downstream human activities remain important bottlenecks.
Similarly, Demirer et al. (2026a) find that AI generates the largest productivity gains when it can execute
contiguous chains of activities, highlighting the importance of workflow structure and task adjacency
for automation and job redesign.
Recent workplace evidence reinforces these arguments. Dillon et al. (2025) study a field experi-
ment in which workers received access to an integrated generative AI tool for email, meetings, and
writing. They find meaningful individual-level time savings, especially in email, but limited evidence
that access changed the quantity or composition of workers’ tasks. Similarly, Chen et al. (2026) show
that the labor-market effects of generative AI depend on whether occupations are more exposed to au-
tomation or augmentation. Using job-posting data, they find reduced demand and skill requirements in
automation-prone jobs but increased demand and skill complexity in jobs involving human–AI collabo-
ration. Together, these findings suggest that realizing the value of AI depends not only on technological
capability but also on how organizations redesign workflows, responsibilities, and skill requirements
around new forms of human–AI collaboration.
Our evidence on differences across Individual, Organizational, and OpenAI users speaks directly
to these organizational-complement mechanisms. If agentic AI adoption depended only on model ca-
pability, we would expect similar patterns of use wherever the same model is available. Large differ-
5
ences across user groups instead suggest that adoption depends on context: access to relevant files and
systems, management expectations, workforce skills, and the availability of complementary review pro-
cesses. In this sense, agentic AI is not simply a cheaper input into existing work; its value depends on
whether organizations can redesign workflows, responsibilities, and review processes around delega-
tion and verification.
3
Who Uses Agentic AI?
We begin by studying who uses Codex, and how adoption of agentic tooling differs across user popu-
lations and worker groups. Codex usage grew rapidly after launch, with the number of weekly active
users increasing more than fivefold between January 1 and June 1, 2026. This establishes the aggregate
growth of Codex, but user counts alone do not capture the shift from conversational to agentic tooling.
Many users continue to use ChatGPT, and users who adopt Codex often use it much more intensively
than users who remain only on conversational interfaces. For that reason, after documenting aggregate
weekly active usage, we focus primarily on Codex’s share of output tokens across Codex and ChatGPT.
Figure 1 compares the shift to Codex across OpenAI workers, organizational users, and individual
users. Panel A shows the extensive margin: whether active users of either product use Codex at all.
Panel B shows the intensive margin: the share of output tokens produced through Codex rather than
ChatGPT. Among individual users, conversational interfaces continue to dominate: fewer than 1\% of
active individual users used Codex in the last 28 days. But the individual users who adopt Codex are
unusually intensive, so Codex accounts for a much larger share of individual-user output tokens than
of individual active users. Among organizational users, adoption is broader. In the last 28 days, 17.3\%
of organizational users used Codex, several times the corresponding share among individual users, and
Codex accounts for the majority of output tokens generated by organizational users in recent weeks. The
most pronounced shift is among OpenAI workers. Codex use is both pervasive and intensive: almost all
active OpenAI workers use Codex each week, and in recent weeks Codex accounts for more than 99\%
of output tokens generated across Codex and ChatGPT.
To assess whether Codex’s growth is driven mainly by developers or by a broader set of users,
we classify Codex users into three broad personas based on their recent Codex requests: Developers,
General Knowledge Workers, and Personal users. Each request receives a persona label from an automated
classifier, and each user is assigned the persona that appears most often in their Codex activity over
the prior 30 days, with ties broken by the most recent request. The prompt used for this classification is
provided in Appendix D.7
Software development activity is assigned the Developer label, other work activity is labeled as Gen-
eral Knowledge Worker, and non-work usage is combined into the Personal bucket. Typical Developer
use includes writing, reviewing, or refactoring code. General Knowledge Worker use includes writing
7We validated the persona classifier using a small sample of employees. For each employee, we compared the assigned persona
against their HR job title to assess alignment. The results showed strong agreement: over 90\% of engineers were classified into
the Developer persona, while more than 90\% of sales employees were classified into the General Knowledge Worker persona.
6
Fig. 1: Codex usage relative to ChatGPT, by user population
(a) Share of active users using Codex
(b) Share of output tokens from Codex
Panel A shows the share of users active on either ChatGPT or Codex during the preceding twenty-eight days who were active
on Codex at least once during that period. Panel B shows output tokens produced on Codex as a share of output tokens
produced on either Codex or ChatGPT during the preceding twenty-eight days. We compare individual users, organizational
users, and OpenAI workers.
or editing documents, analyzing spreadsheets, drafting memos, coordinating communication, and other
non-coding tasks. Personal users tend to use Codex for tasks including hobbies, financial activities, and
education.
Figure 2 shows that Codex adoption is not limited to the initial developer base. Developers remain
an important share of users, especially among individual and organizational accounts, but growth is
faster among Non-developers. Appendix Figure A7 provides a more detailed decomposition of personas
7
Fig. 2: Growth in active Codex users by account type and persona
Figure shows growth in users active in the previous 28 days, by persona and account type. The Non-developer series pools
users assigned to the Personal and General Knowledge Worker personas. Personas are assigned using the most common label
assigned to a user’s queries in the 30 days prior to June 11, 2026. Users’ account types are allowed to vary over time. All series
are indexed to have a value of 1 as of August 1, 2025. Estimates are based on a 3\% sample of users.
for individual and organizational users. It shows that, in contrast to use within OpenAI, individual and
organizational use remains more concentrated in developer personas, especially full-stack engineering.
We next study adoption by job title and function, moving from behavior-based personas to orga-
nizational role measures. For organizational users outside OpenAI, we focus on the subset of firms for
which we have high-quality job-title information.8 We classify job titles into cleaned job-title groups,
seniority levels, and people manager statuses using the classifier described in Appendix C. For OpenAI
workers in this analysis, we instead use directly observed internal job-function metadata.9 These mea-
sures allow us to study whether the shift from ChatGPT to Codex is concentrated in technical roles or
whether it extends across the organization.
Figure 3 highlights broad patterns of adoption. Across all job functions, the share of output tokens
generated on Codex is increasing, with the most rapid progress among technical job functions. Among
Organizational account users, the average engineer produces 26.8\% of their tokens on Codex, a share
that has quintupled since the start of the year. Other technical roles, such as Data and Analytics Practi-
tioners, have seen similarly rapid increases, with the average user in such a role now generating about
15.2\% of their output tokens on Codex. Though non-technical roles have seen a large proportional in-
8Within this set of firms, job-title coverage is broad: for the median firm, we observe job titles for 84.2\% of workers.
9We use classifier-based measures for OpenAI users elsewhere in the paper. In this analysis, however, internal job-function
metadata provides a higher-quality role measure, so we use it instead of classifier-derived job-title groups.
8
Fig. 3: Codex share of output tokens for the average user, by worker type
(a) Organizational users, by inferred job title
(b) OpenAI workers, by job function
Panel A shows the average user-level share of output tokens across ChatGPT and Codex that were produced on Codex among
organizational users in firms with high-quality job-title information, by inferred job-title class. Panel B shows the corresponding
employee-level share among OpenAI workers, by job function. Figures are based on output tokens generated on each platform
in the preceding 28 days. Only users active in the preceding 28 days are included.
9
crease, the average user in a non-technical role generates a relatively small share of their tokens in Codex,
with Codex usage among legal teams accounting for only 1.9\% of output tokens for the average user.
Among Organizational account users, the relatively low share of tokens generated across job func-
tions is primarily due to the limited share of users who use Codex at all, which we highlight in Figure 1.
Despite the fact that Codex usage has not diffused fully in most job functions, it nevertheless accounts
for a large share of overall tokens generated among many job functions, since usage intensity is much
higher among users who adopt Codex. For instance, although Codex accounts for just 26.8\% of the out-
put tokens generated by the average engineer with an Organizational account, Codex usage accounts for
88.3\% of the total output tokens generated by users of this type. This reflects both its adoption among
more intensive users within each function and the deeper use cases agentic tooling enables. We see a
similar divide in non-technical roles: among legal users, Codex accounts for 17.6\% of total output to-
kens, despite accounting for only 1.9\% of tokens generated by the average user.
Usage within OpenAI provides a glimpse into a similar transition where workers began to shift
usage to Codex from conversational harnesses much earlier. In the second half of 2025, usage patterns
within OpenAI looked quite similar to those in external organizational accounts today; engineers used
agentic harnesses for the largest share of their usage, while usage in other job functions lagged behind.
Over the following months, as tooling improved and Codex usage diffused within the company, Codex
continued to rise in all departments, with usage accounting for more than 90\% of the average engineer’s
usage by March 2026, and for more than 90\% of use in almost all departments by the present day.
The transition to agentic tooling was more rapid in later-adopting job functions within OpenAI
than it was among the early-adopting technical functions. In January 2026, when engineering users had
already shifted more than half of their AI usage to Codex, use in functions such as legal and recruiting
was close to zero. From that point, usage increased gradually to about 20\% of output tokens by early
April, then rose quickly from roughly 20\% to 75\% within a month. Although differences across functions
remain, convergence toward agentic tooling was rapid: between December 2025 and April 2026, OpenAI
moved from a pattern in which most functions primarily used conversational AI to one in which Codex
was dominant across functions. This period coincided with a broader internal effort to adopt Codex
across workflows, including training sessions and regular feedback loops.
Finally, Figure 4 shows that adoption among Organizational users is not confined to a single senior-
ity level. Codex’s share of output tokens rises across the seniority distribution, indicating that agentic
tooling is used by both more junior and more senior workers. This result matters because Codex is
not only a tool for direct implementation by individual contributors. Senior users may also use agentic
tooling for planning, review, and other work that involves delegating tasks and evaluating outputs. We
return to these differences in delegated work in Section 4.
The evidence in this section shows where the shift to Codex is occurring: adoption is broadest and
most intensive among organizational users and OpenAI workers, it begins among technical workers
before spreading to non-developer roles, and is occurring across all seniority levels within organizations.
10
Fig. 4: Codex share of output tokens among Organizational users, by inferred seniority level
Each series shows the average user-level share of output tokens across ChatGPT and Codex that were produced on Codex
for organizational users with high-quality job-title information, by inferred seniority level. Figures are based on output tokens
generated on each platform in the preceding 28 days. Only users active in the preceding 28 days are included.
The next section asks what these users do with Codex once they adopt it. That task-level evidence is
important because account type, persona, and job function do not by themselves reveal whether Codex
is being used for software development, broader knowledge work, or personal tasks.
4
What kinds of work do people do with Agentic AI?
This section studies what users delegate to Codex once they adopt it. We use automated classifiers to
measure two aspects of delegated work: the type of task being attempted and, for a smaller sample,
the estimated human time required to complete the task without AI. We first compare task composition
across individual users, organizational users, and OpenAI workers. We then examine whether users are
delegating more complex tasks over time. Finally, we compare task mix across personas, job-title groups,
and seniority levels to show how Codex use differs across worker types.
4.1
Task shares for Codex usage
To measure what users delegate to Codex, we classify Codex requests into a fixed two-level task tax-
onomy. The classifier assigns each request to a primary task category based on the user’s requested
outcome, rather than on incidental tool use during execution. The top-level categories include software-
related work, such as code implementation, code understanding, code validation, engineering opera-
tions, and application management, as well as broader knowledge-work categories such as data analy-
11
sis, research, knowledge artifacts, collaboration, and business-function workflows. The full prompt and
label definitions are reported in Appendix E. We use this taxonomy to compare the tasks Individual,
Organizational, and OpenAI users delegate to Codex.
In Figure 5, we compare the composition of top-level task usage across account categories, measur-
ing each user’s distribution of classified Codex requests and then averaging within account category.10
Across all account types, the most common task categories are tied to software production. Engineering
operations, code implementation, code understanding, application management, and code validation
account for a large share of Codex activity across groups. This indicates that Codex use, while concen-
trated in software production, is not limited to producing new code. Users also rely on it to understand
existing code, validate code changes, manage applications, and support engineering workflows. Impor-
tantly, this suggests that Codex is integrated into the broader software development life cycle rather
than used only as a code-generation tool.
At the same time, firms closer to the frontier of Codex adoption also use AI for more experimental
purposes. Organizational users are more concentrated in coding activities, while OpenAI users show
more activity in research, business function workflows, and engineering operations. The pattern sug-
gests that OpenAI use may involve more experimental, research-oriented, or infrastructure-adjacent
activity, whereas organizational use is more centered on direct software implementation and under-
standing.
4.2
Task complexity on Codex
The scope of possible requests within each task category is broad: the Code Implementation category
covers queries requesting one-line tweaks as well as queries requesting the creation of full-fledged soft-
ware applications. We quantify the complexity of users’ use cases using a prompt, which analyzes the
text of a subset of queries and estimates the time that it would take for an experienced human worker to
complete the task without the assistance of AI. We use the classifier to assign estimated task durations to
the prompts sent by a 0.1\% random sample of Individual accounts who have opted to allow their data to
be used for model training. By observing the estimated completion time of turns since November 2025,
we are able to measure how the complexity of tasks performed by Codex has changed, as the underlying
models have improved and humans have gained experience in assigning tasks to Codex. We provide
the prompt used, as well as validation metrics, in Appendix F.
In Figure 6, we present statistics about the distribution of query complexity. Over time, the complex-
ity of users’ requests has increased substantially—in December 2025, only 35.4\% of active Individual
users sent at least one prompt that we estimate would have taken an experienced human at least one
hour to complete. This share has increased rapidly—in May 2026, we estimate that 70.2\% of users sent
at least one prompt of similar or greater difficulty. We estimate even larger proportional increases in the
frequency at which users send more difficult prompts: since December 2025, the share of users who send
10Appendix B discusses complementary ways of analyzing task usage, including measuring the share of users who use each
category rather than the average user’s task mix, and decomposing task use into more granular “second-level” categories.
12
Fig. 5: Share of Codex usage by first-task for the average user
Shares reflect the average user-level share of classified messages falling into each first-level task category during the 28 days
prior to June 11, 2026. The prompt used for our classifications is provided in Appendix E Messages classified as “other” are
excluded from the calculation. Organizational and Individual calculations are based on a 4\% sample of users; users inactive
during the period are excluded. Error bars show 95\% cluster bootstrap intervals from resampling sampled Individual and
Organizational users; intervals reflect sampling uncertainty.
at least one prompt that would have taken an experienced human 8 hours to complete has risen from
2.1\% to 25.6\%.
The most complicated queries are concentrated at the start of users’ sessions. In Figure 6b, we show
that the first turn of a thread is more than twice as likely as the fourth turn to require action that would
take an experienced human more than one hour. This suggests that users tend to request the broadest,
most complex portion of work initially, with further, narrower refinements following later in the session.
4.3
Task shares by worker type
We next ask whether the task mix differs across the worker groups studied in Section 3. Combining the
task classifier introduced above with job title and seniority measures allows us to compare the com-
position of Codex use within each group. Because the figures report within-group token shares, they
describe differences in what users delegate to Codex rather than differences in overall adoption or us-
age intensity.
Figure 7 compares the composition of Codex use across job title groups, separately for OpenAI
13
Fig. 6: Model-estimated complexity of Codex turns among Individual account users
(a) Distribution of users by peak task complexity
(b) Complexity of turns by turn index
Durations reflect the time it would take an experienced human to complete each task, estimated using the classifier described
in Section F. Each panel is constructed using a sample of 0.1\% of Individual users who have opted to allow their queries to be
used in training. Figure 6a presents the share of sampled users each month that ask at least one query over each threshold of
estimated human completion time. The data for February and March is partial. Figure 6b shows the share of turns estimated
to have a human completion time over certain thresholds, split by the query’s position in the thread. Only messages sent in
May 2026 by Individual users in the 0.1\% sample are included. An index of 1 corresponds to the first message sent in a thread.
workers and users in organizations with job title coverage. We see that most usage both inside OpenAI
and inside organizations more broadly is concentrated in engineering operations. Inside OpenAI, across
developer and non-developer roles, knowledge artifacts, collaboration, and application management
are common tasks. In organizations, many workers are using agentic tools for generating knowledge
artifacts, especially in sales, marketing, and recruiting functions.
Figure 8 compares the distribution of Codex output tokens across task categories by inferred worker
seniority, separately for OpenAI employees and for workers within organizations with job title cover-
age. We find that Codex use differs meaningfully across seniority groups, but the patterns do not sup-
port a simple interpretation in which junior workers use Codex primarily for code generation while
senior workers use it primarily for managerial tasks. In organizations, engineering operations accounts
for the largest share of output tokens for most seniority groups, while code implementation is espe-
cially prominent among early-career workers and contractors. Knowledge artifacts account for a larger
share of tokens among executives and managers, suggesting that more senior workers use Codex partly
to produce, interpret, and organize technical documentation and other written outputs. Within Ope-
nAI, the distribution is more uneven: Individual Contributors (ICs) and senior ICs allocate relatively
large shares of tokens to engineering operations, while managers and executives devote larger shares to
collaboration-oriented tasks. Overall, the figure shows that workers across the seniority distribution use
Codex for a broad portfolio of technical and organizational activities, including engineering operations,
implementation, code understanding, validation, research, data analysis, and documentation.
14
Fig. 7: Differences in Codex use cases across different types of workers, both internally at
OpenAI (top) and in organizations with job title coverage (bottom)
Distribution of Codex tokens across task types, by inferred job title, within OpenAI (top) and organizations with job title
coverage (bottom).
5
How Are People Using Agentic AI?
Although agentic AI can already perform tasks within existing production processes, the literature on
general-purpose technologies suggests that the largest productivity gains often arise when firms reorga-
nize production around the new technology rather than merely substitute it into existing workflows. We
draw here on the canonical example in David (1990), who studies the delayed productivity gains from
the transition from steam power to electric power in manufacturing. In the early stages of electrification,
many factories replaced centralized steam engines with centralized electric motors while preserving ex-
isting factory layouts and work patterns. The larger gains emerged only when firms redesigned produc-
15
Fig. 8: Differences in Codex use cases by seniority level and account type
Distribution of Codex tokens across task types, by inferred seniority level, within OpenAI (top) and organizations with job
title coverage (bottom).
tion around the distinctive capabilities of electric power: smaller motors made it possible to decentralize
power, reorganize the factory floor, change the sequencing of tasks, and adopt more flexible production
layouts. These changes took decades to diffuse, partly because firms had to incur substantial fixed costs
to redesign plants and develop complementary organizational practices. The history of technological
adoption suggests that the near-term effects of agentic AI may understate its longer-run productivity
potential if firms have not yet discovered, adopted, or scaled the new production processes that the
technology makes feasible.
Building on this literature, we next examine whether users are beginning to adapt their workflows
to the capabilities of agentic systems. Having shown what tasks users delegate to Codex, we now study
how that delegation is organized. We focus on three observable margins: whether users run multiple
agents in parallel, whether agents perform longer blocks of work on users’ behalf, and whether users
16
Fig. 9: Codex turn concurrency
Figure summarizes concurrent Codex usage in the week prior to June 11, 2026, plotting the distribution, across account types, of
users’ peak number of concurrent agents. Simultaneous agents are measured using overlapping turns, spread across different
threads, which overlap for more than 30 seconds
systematize repeated work through skills and plugins. These measures do not fully capture organiza-
tional redesign, but they provide early evidence on whether agentic AI use is moving beyond one-off
assistance toward more persistent, repeatable, and parallel workflows. The speed of this adjustment
may reflect an important difference between digital production processes and earlier episodes of tech-
nological change. Whereas electrification often required firms to redesign physical plants and replace
durable capital, agentic AI allows workers and organizations to experiment with new workflows at
relatively low cost. This lower cost of experimentation may allow new production methods to diffuse
more quickly than in prior general-purpose technology transitions, even though the full organizational
complements to agentic AI are still emerging.
5.1
Turn Concurrency
Codex, like many AI agents, uses a threaded interaction model in which users can initiate multiple
agents and interact with each one in a largely independent workspace. This design changes the structure
of delegation. Because long-running tasks may require an agent to work independently for minutes or
hours, users need not wait for one task to finish before beginning another. Instead, they can assign
several tasks to different agents, monitor progress across threads, and intervene selectively when an
agent requires clarification, review, or correction. The threaded model therefore allows users to manage
multiple streams of agentic work in parallel, expanding the scope for concurrent task execution relative
to conventional single-threaded interactions with AI systems.
It has been speculated that, as agentic tooling advances, parallel workflows will become more com-
mon, increasing the value of skills such as delegation and management (Alekseeva et al. 2026). In Fig-
ure 9, we explore the degree to which users are currently using Codex in parallel fashion. For each user,
we calculate the number of overlapping turns they have in different threads during the week prior to
17
June 11, 2026.11 Each bar shows the distribution of peak concurrency during that week, among users of
each account type. Among Organizational and Individual users, the usage of concurrent turns is fairly
minimal. Roughly 67.4\% of Organizational users and 63.9\% of Individual users do not use concurrent
turns at all during this period, and among those who do, the majority peak at two concurrent turns.
Among OpenAI users, however, workflows have become far more parallelized. Only 10.7\% of users in
this group use a sole workflow at any one time during this period, and nearly 28.6\% managed five or
more concurrent agents at some point during this period. The prevalence of this practice within OpenAI
represents a workflow in which humans oversee a team of agents, delegating tasks across many si-
multaneous workers at once. This workflow is fundamentally different from the one practiced by most
external users, in which humans directly perform the majority of work, as it requires workers to manage,
delegate to, and review the work of a relatively large group of agents.
5.2
Long-running agents
Parallelism captures one way agentic use departs from conversational interaction; another is the dura-
tion of delegated work. Researchers and practitioners have emphasized the possibility that agents can
operate asynchronously for extended periods with limited human supervision.
In Figure 10a, we examine whether this possibility has begun to reshape workflows within OpenAI
by measuring the amount of time Codex tasks remain active over the course of a day. For the median
OpenAI employee, agentic workflows remain intermittent: on June 11, 2026, the median employee had
Codex turns running for 2.5 hours. The typical user therefore delegates meaningful blocks of work, but
has not shifted to continuous, around-the-clock autonomous execution.12 Among frontier users, how-
ever, this pattern looks considerably different; OpenAI employees at the 99th percentile of the distri-
bution have recently run about 71 hours of agent turns within the average day, which implies that at
a given hour of the day, they have several agents running concurrently. These parallel, long-running
workflows have increased in prevalence recently, with 99th percentile users within OpenAI increasing
their cumulative daily runtime by nearly 88\% since April 7, 2026. During the period we observe, the
absolute increase in daily runtimes within OpenAI was largest at the top of the distribution, although
lower percentiles grew faster proportionally.
Outside OpenAI, cumulative daily runtime remains substantially lower, and the median user con-
tinues to use Codex only intermittently over the course of the day. We nevertheless observe substantial
growth at the upper tail of the usage distribution during our sample period. At the 99th percentile, daily
runtime increased by roughly 25 percent among Organizational users and by about 50 percent among
Individual users. While usage intensity is rising across all user populations, these patterns show that
OpenAI employees use Codex much more intensively than external users. The contrast between the
11We restrict only to turns which run in different threads, and which overlap one another for at least 30 seconds.
12We measure the time an agent is active within each turn by summing the latency of all completed request-response pairs. We
include gaps of up to 30 minutes between responses, which include the results of tool calls, reasoning, and the output text
the agent returns. We remove longer gaps, which are often driven by agents who are awaiting user input or paused. Each
user is assigned the total duration of all agent turns they initiate and which begin within that calendar day. Turns that are
not interactive, such as those that are run as part of a batch script, are excluded.
18
Fig. 10: Cumulative Codex turn duration per day
(a) OpenAI users
(b) Organizational users
(c) Individual users
Figure summarizes total active runtime of users’ Codex turns on each day. This measure excludes gaps where no output is
returned for more than 30 minutes, which predominantly reflect periods during which Codex is awaiting user input. Because
overlapping turns are summed, a user’s cumulative runtime during a day can exceed 24 hours. Figure 10a plots usage among
OpenAI accounts, Figure 10b plots usage among Organizational accounts. Figure 10c plots usage among Individual accounts.
Accounts inactive on each day are excluded.
median and the upper tail suggests that agentic workflows remain sporadic for typical users, but that a
smaller group of high-intensity users is rapidly expanding the amount of work it delegates to Codex.
5.3
Systematization of Agentic Work
A final margin is whether delegated workflows remain one-off interactions or become reusable routines.
The simplest use of agentic AI is ad hoc: a user describes a task, the agent performs it, and the interaction
ends. Codex, like other agentic AI software, also allows users to codify workflows and procedural pref-
erences so that similar work can be delegated repeatedly. We refer to this shift from ad hoc use to reusable
workflow infrastructure as the “systematization” of agentic work. Systematization is an important step
toward delegated production: without it, users must repeatedly supply task context, procedural guid-
ance, and instructions, limiting the extent to which work can be handed off.
19
Fig. 11: Skill use over time and by account category
(a) Skill use over time
(b) Skill use by account category
(a) Each series shows the share of weekly active Codex users who invoked each skill source at least once during
the prior 7-day window. Users can appear in multiple skill-source categories. (b) Each dot shows the share of
weekly active Codex users in an account category who invoked each skill source at least once during the 7-day
window ending on June 11, 2026. Users can appear in multiple categories.
20
Within Codex, work is often systematized through skills and plugins, which combine instructions,
software, and external-tool integrations that can be reused during a task.13 These skills and plugins can
be shared across users, allowing workflows to be systematized not only by individuals but also within
and across organizations. They therefore make it easier to transfer procedural context, tool access, and
task-specific guidance across repeated uses of Codex.
In our analysis, we distinguish between five sources of skills:
• Preinstalled skills, such as image generation, are reusable capabilities bundled with Codex.
• Curated skills, such as PDF handling, are standalone OpenAI-distributed skills that are not tied to
a plugin.
• Plugin skills, such as Google Drive document workflows, are bundled into a plugin alongside other
skills or software.
• Custom plugin skills are skill invocations that are associated with a recognized plugin but do not
match a skill in the curated plugin skill catalog.
• Custom skills are standalone skills not distributed by OpenAI, such as team-specific data visual-
ization guidelines or a research workflow.
These categories capture different forms of systematization, ranging from product-provided reusable
capabilities to user- or organization-specific workflow codification.
Figure 11 shows that skill use is already common, increasing over time, and unevenly distributed
across account categories. In the 7-day window ending on June 11, 2026, 25.7\% of active Individual
Codex users and 30.4\% of active Organizational Codex users invoked at least one skill. Within OpenAI,
skill use is nearly universal: 96.2\% of active Codex users invoked at least one skill. Skill use also increased
substantially during the period shown in Figure 11a: the share of active Codex users invoking any skill
rose from 5.4\% on March 1, 2026 to 26.6\% on June 11, 2026.
The growth in skill use comes especially from plugins and custom skills. These two categories repre-
sent different forms of systematization. Plugins extend Codex’s general-purpose capabilities into recur-
ring task domains such as documents, spreadsheets, slide decks, and other structured artifacts. Custom
skills serve a different role: they allow users and organizations to encode more local procedural con-
text, such as team writing standards, recurring reports, organization-specific workflows, or user-specific
preferences. The growing use of custom skills suggests that users value not only the model’s general
capabilities, but also the ability to attach persistent procedural context to repeated tasks that are too
specialized to be fully covered by standardized plugins.
13In OpenAI’s terminology, a skill is the authoring format for a reusable, task-specific workflow: a directory containing a
required SKILL.md file, with name and description metadata and optional scripts, references, and assets. A plugin is
an installable distribution unit, identified by a .codex-plugin/plugin.json manifest, that can package skills together
with app integrations, Model Context Protocol (MCP) configuration, hooks, and supporting assets. Thus, a skill specifies a
workflow, whereas a plugin packages and distributes capabilities.
21
Skill use differs across account types not only in terms of overall frequency, but also in terms of the
composition of skills invoked. Figure 11b shows that plugin and custom skill usage is highest among
OpenAI users and substantially higher among Organizational users than among Individual users. These
differences are consistent with variation in the returns to workflow systematization across settings. Cus-
tom skills are especially valuable when repeated tasks depend on organization-specific context, shared
conventions, internal procedures, or team-level standards. These conditions are more common in Or-
ganizational settings than in Individual use. The large differences in custom skill adoption therefore
suggest that users derive the greatest value from systematized workflows in high-context organiza-
tional environments, where persistent procedural context can reduce coordination costs and improve
consistency across repeated tasks.14
6
Conclusion
We document the emergence of agentic AI as a distinct workplace technology. Much prior research
has studied AI as an information, advice, and content-generation tool, where the primary interaction
is conversational. Agentic systems introduce a different mode of use. Rather than simply asking for
information, users increasingly delegate work: producing artifacts, modifying systems, and executing
workflows. This shift suggests that standard measures of AI use, such as active users, chats, or message
volume, may become less informative as agentic systems diffuse. Future analyses may need to track
delegated task complexity, runtime, workflow reuse, concurrency, and production output.
The evidence we report suggests that the diffusion of agentic AI is not just a shift in which AI tool
people use, but a shift in how they organize work around AI. Early use often resembles familiar con-
versational assistance: asking questions, drafting text, or generating code in a back-and-forth exchange.
More agentic use involves assigning tasks, reviewing outputs, coordinating multiple threads of work,
and reusing codified workflows. Intensive users appear to manage portfolios of agentic work, with
their own role shifting toward delegation, supervision, and integration. This suggests that the adoption
of agentic AI involves the development of new work practices, not simply more frequent use of AI.
Software development is the leading edge of this change. It is still the most common use case for
agentic AI tools across user populations, likely because software work has properties that lend them-
selves to agentic AI: software work is digital, produces verifiable artifacts, and contains many modular
subtasks. At the same time, observed uses extend beyond code generation. Users are also delegating
system understanding, debugging, validation, documentation, configuration, and operational work.
We find that as adoption deepens, the scope of agentic work expands. The broadest task portfolios
are observed where frictions to adoption are lowest. Among OpenAI workers, Codex is used not only for
software engineering but also for research, planning, communication, recruiting, sales, product work,
and data analysis. These patterns suggest that while software may be the initial deployment domain for
14Appendix Figure A8 shows that skill invocation rates also vary by task area. Among Individual and Organizational users,
skill use remains concentrated in technical workflows, including application management, collaboration, code validation,
and code understanding. Within OpenAI, skill use has diffused further into broader knowledge-work workflows; for exam-
ple, 50.9\% of collaboration conversations invoke at least one skill.
22
agentic AI, it is diffusing quickly across this boundary. Figure 12 shows that as usage patterns within
OpenAI have intensified, so have measures of output—between November 1, 2025 and June 11, 2026,
the median worker’s number of output tokens rose at least 10-fold in every job function. This growth in
token output reflects several factors, including the shift from conversational interfaces to agentic tools
and the later emergence of more intensive, persistent uses of those tools. Our evidence from OpenAI is
not necessarily representative of the typical organization, but it shows what agentic AI use can look like
when adoption frictions are low: use expands across functions, output rises sharply, and work increas-
ingly takes the form of delegated tasks rather than isolated exchanges.
Fig. 12: Change in median output tokens per person, by OpenAI job function
Plot shows the normalized number of output tokens generated in the trailing 28 days by the median worker in each job function
within OpenAI. Output tokens include those generated on ChatGPT and those generated on Codex. Each point includes only
users active in the preceding 28 days. All series are normalized to have a value of 1 on November 1, 2025.
These patterns have implications beyond adoption and productivity measurement. As AI use shifts
from consultation to delegation, organizations may need to rethink how work is allocated, reviewed,
and coordinated. Agentic systems make it possible to offload larger, more repeatable, and more modu-
lar units of work, potentially reducing the importance of some routine execution tasks while increasing
the importance of judgment, oversight, coordination, and review. Jobs may increasingly involve direct-
ing, monitoring, and integrating the outputs of AI agents rather than executing each component task
directly. Organizations may also redesign workflows around parallelized agent labor, allowing individ-
ual workers to manage multiple workstreams at once. These changes could affect team composition,
hiring needs, career ladders, and the distribution of work across skill levels.
Taken together, the evidence suggests that agentic AI is not simply a more capable form of con-
23
versational AI. Conversational systems are mainly used to exchange information with users; agentic
systems are increasingly used to carry out work on users’ behalf. The most intensive users do not just
chat more often. They delegate tasks, run work in parallel, reuse codified workflows, and shift their own
effort toward supervision and integration. The frontier of AI adoption is therefore moving from asking
systems for answers toward managing systems that act. Documenting that transition is essential for
understanding how AI may reshape work, organizations, and the structure of knowledge production.
References
Alekseeva, L., Azar, J., Giné, M., and Samila, S. (2026). “Artificial Intelligence Adoption and the Demand for
Managerial Expertise”. Strategic Management Journal. Online first, 1–27. DOI: 10.1002/smj.70099.
Baumann, J., Padmakumar, V., Li, X., Yang, J., Yang, D., and Koyejo, S. (2026). “SWE-chat: Coding Agent Interac-
tions from Real Users in the Wild”. arXiv preprint arXiv:2604.20779. DOI: 10.48550/arXiv.2604.20779.
Bick, A., Blandin, A., and Deming, D. J. (2026). “The Rapid Adoption of Generative AI”. Management Science.
Advance online publication. DOI: 10.1287/mnsc.2025.02523.
Brynjolfsson, E., Rock, D., and Syverson, C. (2019). “Artificial Intelligence and the Modern Productivity Paradox:
A Clash of Expectations and Statistics”. The Economics of Artificial Intelligence: An Agenda. 
\end{document}
